% Artigo IEEE (coluna dupla) — Estudo comparativo CNN vs ViT na detecção de imagens geradas por IA.
% Números e figuras vêm de results/ via `python -m src.export_paper` (pastas data/ e figures/).
\documentclass[conference]{IEEEtran}
\IEEEoverridecommandlockouts

\usepackage[T1]{fontenc}
\usepackage[utf8]{inputenc}
\usepackage[brazil]{babel}
\usepackage{cite}
\usepackage{amsmath,amssymb}
\usepackage{graphicx}
\usepackage{booktabs}
\usepackage{xcolor}
\usepackage{url}
\usepackage[hidelinks]{hyperref}
\graphicspath{{figures/}}

% Macros numéricas geradas pelos scripts (src/eda.py, src/stats.py, src/xai.py).
\input{data/eda_macros.tex}
\input{data/stats_macros.tex}
\InputIfFileExists{data/xai_macros.tex}{}{%
  \newcommand{\XaiN}{--}\newcommand{\DelCnnCam}{--}\newcommand{\DelCnnRnd}{--}%
  \newcommand{\DelVitCam}{--}\newcommand{\DelVitRnd}{--}}

\begin{document}

\title{CNN ou Vision Transformer? Um Estudo Comparativo com 30 Simulações Pareadas
na Detecção de Imagens Geradas por Inteligência Artificial}

% TODO: preencher autores e instituição.
\author{\IEEEauthorblockN{Autor 1, Autor 2, Autor 3}
\IEEEauthorblockA{Programa / Disciplina\\
Instituição — Cidade, Brasil\\
\{autor1, autor2\}@instituicao.edu.br}}

\maketitle

\begin{abstract}
A popularização de geradores de imagem torna necessário distinguir automaticamente
imagens reais de imagens sintéticas. Este trabalho compara duas abordagens que
compartilham \emph{Transfer Learning} e explicabilidade (Grad-CAM) e diferem apenas
no classificador: uma rede convolucional (ResNet-50) e um \emph{Vision Transformer}
(ViT-S/16), com número de parâmetros equivalente. Sobre um conjunto público de
\NImages{} imagens (\NReal{} reais e \NAi{} geradas por IA), foram executadas
\NSim{} simulações pareadas com partições estratificadas distintas. A CNN obteve
F1-macro de \FoneCnnMean{}~$\pm$~\FoneCnnStd{} e o ViT, \FoneVitMean{}~$\pm$~\FoneVitStd{}.
O teste \FoneTest{} sobre as diferenças pareadas resultou em $p=\FonePvalue{}$, de modo
que a hipótese nula de desempenho igual foi \Rejected{} ao nível de 5\%.
Mapas Grad-CAM e a métrica de \emph{deletion} indicam em que regiões cada modelo
baseia suas decisões.
\end{abstract}

\begin{IEEEkeywords}
detecção de imagens sintéticas, transfer learning, vision transformer, CNN,
Grad-CAM, teste de hipótese
\end{IEEEkeywords}

% =====================================================================
\section{Introdução}
Modelos generativos de difusão e GANs produzem hoje imagens que observadores humanos
têm dificuldade em distinguir de fotografias e obras reais, com impacto em
desinformação, direitos autorais e confiança em mídias digitais. Detectores
baseados em aprendizado profundo são a principal resposta técnica a esse
problema~\cite{wang2020cnnspot,corvi2023diffusion,ojha2023universal}.

Uma pergunta recorrente é qual família de arquitetura é mais adequada: as redes
convolucionais (CNNs), com viés indutivo de localidade e sensíveis a texturas e
artefatos de alta frequência, ou os \emph{Vision Transformers} (ViTs), que modelam
relações globais por atenção~\cite{dosovitskiy2021vit,raghu2021vitcnn}. Comparações
de uma única execução, porém, são frágeis: com conjuntos pequenos, a variação causada
pela partição dos dados pode ser maior que a diferença entre modelos.

Este trabalho conduz um estudo comparativo controlado com três técnicas vistas na
disciplina: (i)~\emph{Transfer Learning}, (ii)~classificação por CNN \emph{vs.} por
Transformer (técnica contrafactual) e (iii)~explicabilidade (xAI) com Grad-CAM.
As contribuições são: um protocolo de \NSim{} simulações pareadas com teste de
hipótese formal; a comparação de desempenho, custo computacional e explicações
entre os dois cenários; e o código reprodutível.

% =====================================================================
\section{Trabalhos Relacionados}
Wang \emph{et al.}~\cite{wang2020cnnspot} mostraram que uma ResNet-50 treinada em
imagens de uma única GAN generaliza para outras arquiteturas geradoras quando
treinada com aumento de dados adequado. Corvi \emph{et al.}~\cite{corvi2023diffusion}
estenderam a análise a modelos de difusão, evidenciando traços espectrais
característicos. Ojha \emph{et al.}~\cite{ojha2023universal} usaram o espaço de
\emph{features} de um modelo ViT pré-treinado (CLIP) para obter detectores mais
universais. Bird e Lotfi~\cite{bird2024cifake} propuseram o conjunto CIFAKE e
usaram Grad-CAM para interpretar um classificador convolucional. Raghu \emph{et
al.}~\cite{raghu2021vitcnn} compararam as representações internas de ViTs e CNNs,
mostrando que os ViTs agregam informação global já nas primeiras camadas.

% =====================================================================
\section{Fundamentação}
\subsection{Transfer Learning}
\emph{Transfer Learning} reaproveita um modelo treinado em uma tarefa de origem
com muitos dados (aqui, ImageNet~\cite{deng2009imagenet}) e o ajusta à tarefa
de destino~\cite{pan2010survey}. Com apenas \NImages{} imagens, treinar do zero
seria inviável; por isso ambos os cenários partem de pesos pré-treinados e fazem
\emph{fine-tuning} completo.

\subsection{CNN: ResNet-50}
A ResNet-50~\cite{he2016resnet} empilha blocos convolucionais com conexões
residuais. Utilizou-se a versão \texttt{resnet50.a1\_in1k} da biblioteca
\texttt{timm}~\cite{wightman2019timm,wightman2021rsb}.

\subsection{Transformer: ViT-S/16}
O ViT~\cite{dosovitskiy2021vit} divide a imagem em \emph{patches} de
$16\times16$ pixels, projetados como \emph{tokens} e processados por camadas de
auto-atenção. Utilizou-se \texttt{vit\_small\_patch16\_224.}\allowbreak\texttt{augreg\_in21k\_ft\_in1k}~\cite{steiner2022augreg},
escolhido por ter número de parâmetros próximo ao da ResNet-50 (Tabela~\ref{tab:results}).

\subsection{Explicabilidade: Grad-CAM}
O Grad-CAM~\cite{selvaraju2017gradcam} pondera os mapas de ativação de uma camada
pelos gradientes da classe predita, produzindo um mapa de relevância. Na CNN usa-se
o último bloco residual; no ViT, os \emph{tokens} do último bloco são remodelados
para a grade $14\times14$ de \emph{patches}~\cite{gildenblat2021gradcam}. A fidelidade
das explicações é medida pela métrica de \emph{deletion}~\cite{petsiuk2018rise}:
removem-se progressivamente os pixels mais relevantes e mede-se a queda da
probabilidade da classe predita; quanto menor a área sob a curva, mais fiel a
explicação.

% =====================================================================
\section{Metodologia}
\subsection{Conjunto de dados e pré-processamento}
Utilizou-se o conjunto \emph{AI Generated Images vs Real Images}~\cite{kaggle_dataset},
com imagens de tamanhos e formatos variados. Dos \NFound{} arquivos encontrados,
\NCorrupted{} estavam corrompidos, \NTiny{} tinham lado menor que 32~pixels e
\NDuplicates{} eram duplicatas exatas (mesmo hash MD5), restando \NImages{} imagens
(\NReal{} reais e \NAi{} geradas por IA). A Fig.~\ref{fig:eda} mostra que formato e
resolução originais diferem entre as classes; para impedir que o classificador use
esses atalhos, todas as imagens foram convertidas para RGB, redimensionadas para
menor lado de 256~pixels e salvas em PNG.

\begin{figure}[t]
\centering
\includegraphics[width=\linewidth]{fig_eda_formats.pdf}
\caption{Formato e resolução originais por classe. As diferenças motivam a
padronização das imagens antes do treino.}
\label{fig:eda}
\end{figure}

\subsection{Cenários comparados}
\begin{itemize}
\item \textbf{Cenário A ($H_0$):} Transfer Learning + CNN (ResNet-50) + Grad-CAM.
\item \textbf{Cenário B ($H_1$):} Transfer Learning + ViT-S/16 + Grad-CAM.
\end{itemize}
Apenas o classificador muda entre os cenários (técnica contrafactual).

\subsection{Protocolo de treino}
Os dois modelos usam o mesmo protocolo: entrada $224\times224$; aumento de dados
leve (\emph{RandomResizedCrop} com escala 0,8--1 e espelhamento horizontal);
AdamW~\cite{loshchilov2019adamw} com taxa de aprendizado $10^{-4}$ e
\emph{weight decay} 0,05; agenda cosseno; lote 32; precisão mista; até 10 épocas
com \emph{early stopping} (paciência 3) pelo F1-macro de validação, restaurando os
melhores pesos. A normalização segue a estatística de pré-treino de cada modelo.

\subsection{Simulações e teste de hipótese}
Para cada simulação $s\in\{0,\dots,\NSim-1\}$, o conjunto é particionado de forma
estratificada em 70\%/15\%/15\% (treino/validação/teste) com semente $s$, e ambos os
modelos são treinados e avaliados \emph{na mesma partição}. As amostras são,
portanto, pareadas. A métrica primária é o F1-macro no teste; acurácia e AUC-ROC
são secundárias. Sejam $d_s=F1^{\text{ViT}}_s-F1^{\text{CNN}}_s$ as diferenças:
\begin{equation}
H_0:\ \mu_d = 0 \qquad H_1:\ \mu_d \neq 0, \qquad \alpha=0{,}05.
\end{equation}
A normalidade de $d_s$ é verificada com Shapiro-Wilk~\cite{shapiro1965}; se não for
rejeitada, aplica-se o teste $t$ pareado; caso contrário, o teste de postos
sinalizados de Wilcoxon~\cite{wilcoxon1945,demsar2006}. Reportam-se ainda o intervalo
de confiança de 95\% de $\mu_d$ por \emph{bootstrap}~\cite{efron1993bootstrap},
o tamanho de efeito $d_z$ de Cohen~\cite{cohen1988power} e a correlação
\emph{rank-biserial}~\cite{kerby2014}.

% =====================================================================
\section{Resultados}
\begin{table}[t]
\centering
\caption{Desempenho no teste (média $\pm$ desvio-padrão em \NSim{} simulações) e
custo computacional. $p$: teste principal ($t$ = $t$ pareado, W = Wilcoxon);
$^{*}$: significativo a 5\%.}
\label{tab:results}
\resizebox{\linewidth}{!}{\input{data/table_results.tex}}
\end{table}

A Tabela~\ref{tab:results} e a Fig.~\ref{fig:box} resumem as \NSim{} simulações.
Em F1-macro, o ViT superou a CNN em \FoneVitWins{} simulações e foi superado em
\FoneCnnWins{} (\FoneTies{} empates) (Fig.~\ref{fig:paired}). A diferença média foi
$\bar d=\FoneDiff{}$, com IC~95\% $[\FoneCiLow{};\ \FoneCiHigh{}]$.

\begin{figure*}[t]
\centering
\includegraphics[width=0.9\textwidth]{fig_boxplot_metrics.pdf}
\caption{Distribuição das métricas no teste nas \NSim{} simulações (cada ponto é uma
simulação). O \emph{ensemble} é apenas exploratório.}
\label{fig:box}
\end{figure*}

\begin{figure}[t]
\centering
\includegraphics[width=0.85\linewidth]{fig_paired_f1.pdf}
\caption{F1-macro pareado por simulação: cada linha liga os dois modelos treinados
na mesma partição.}
\label{fig:paired}
\end{figure}

O teste de Shapiro-Wilk sobre $d_s$ resultou em $p=\FoneShapiroP{}$
(Fig.~\ref{fig:diff}), e por isso o teste principal foi o \FoneTest{}. O teste $t$
pareado deu $p=\FoneTP{}$ ($d_z=\FoneDz{}$) e o de Wilcoxon $p=\FoneWilcoxonP{}$
($r_{rb}=\FoneRrb{}$). \textbf{A hipótese nula foi \Rejected{}} ao nível de 5\%.
Para acurácia e AUC-ROC, os valores de $p$ do teste principal foram
\AccPvalue{} e \AucPvalue{}, respectivamente.

\begin{figure*}[t]
\centering
\includegraphics[width=0.9\textwidth]{fig_diff_hist_qq.pdf}
\caption{Diferenças pareadas $d_s$: histograma com IC~95\% \emph{bootstrap} (esq.)
e QQ-plot contra a Normal (dir.).}
\label{fig:diff}
\end{figure*}

As curvas ROC médias (Fig.~\ref{fig:roc}) e as matrizes de confusão agregadas
(Fig.~\ref{fig:cm}) mostram como os erros se distribuem entre as classes. O
\emph{ensemble} exploratório (média das probabilidades) atingiu F1-macro de
\EnsFoneMean{}~$\pm$~\EnsFoneStd{}.

\begin{figure}[t]
\centering
\includegraphics[width=0.85\linewidth]{fig_roc_mean.pdf}
\caption{Curva ROC média $\pm$ desvio-padrão nas \NSim{} simulações.}
\label{fig:roc}
\end{figure}

\begin{figure}[t]
\centering
\includegraphics[width=0.85\linewidth]{fig_confusion_mean.pdf}
\caption{Matrizes de confusão agregadas e normalizadas por linha.}
\label{fig:cm}
\end{figure}

\subsection{Explicabilidade}
A Fig.~\ref{fig:gradcam} mostra mapas Grad-CAM na simulação \XaiSeed{}, escolhida
por ter $d_s$ mais próximo da mediana. Na métrica de \emph{deletion}
(Fig.~\ref{fig:deletion}), calculada sobre \XaiN{} imagens de teste, a área sob a
curva foi \DelCnnCam{} para a CNN (\DelCnnRnd{} com ordem aleatória) e \DelVitCam{}
para o ViT (\DelVitRnd{} com ordem aleatória).

\begin{figure*}[t]
\centering
\includegraphics[width=\textwidth]{fig_gradcam_grid.pdf}
\caption{Grad-CAM: imagem original (topo), CNN (meio) e ViT (base). Colunas
agrupam acertos e erros de cada classe.}
\label{fig:gradcam}
\end{figure*}

\begin{figure}[t]
\centering
\includegraphics[width=0.85\linewidth]{fig_deletion_curves.pdf}
\caption{Curvas de \emph{deletion}: probabilidade da classe predita conforme os
pixels mais relevantes são removidos. Linhas tracejadas: ordem aleatória.}
\label{fig:deletion}
\end{figure}

% =====================================================================
\section{Discussão}
% TODO (após os resultados reais): interpretar o sentido e o tamanho do efeito,
% descrever padrões visuais do Grad-CAM (textura/bordas vs. regiões globais) e
% comparar custo (tempo de treino e inferência, Tabela I).
O protocolo pareado controla a principal fonte de variação em conjuntos pequenos,
a partição dos dados: cada par de modelos é avaliado sobre exatamente as mesmas
imagens de teste, de modo que $d_s$ isola o efeito da arquitetura.

\textbf{Limitações.} O conjunto é pequeno e heterogêneo (imagens de arte, fotografia
e ilustração, de geradores não identificados), e o teste em cada simulação tem
poucas dezenas de imagens, o que aumenta a variância das métricas. As simulações não
são totalmente independentes, pois as partições compartilham imagens. Também não se
avaliou a generalização para geradores não vistos no treino, cenário em que a
literatura reporta quedas acentuadas de desempenho~\cite{ojha2023universal}.

% =====================================================================
\section{Conclusão}
Este trabalho comparou CNN (ResNet-50) e ViT-S/16 na detecção de imagens geradas por
IA, mantendo fixos o \emph{Transfer Learning} e o Grad-CAM. Em \NSim{} simulações
pareadas, a hipótese nula de igualdade de F1-macro foi \Rejected{}
($p=\FonePvalue{}$, \FoneTest{}): \ConclusionText{}.
Trabalhos futuros incluem avaliar generalização entre geradores, conjuntos maiores e
métodos de explicação específicos para Transformers~\cite{abnar2020rollout}.

\section*{Disponibilidade do código}
Código, notebook e resultados: \url{https://github.com/Breno-Lira/Plant-seedlings}.

\bibliographystyle{IEEEtran}
\bibliography{refs}

\end{document}
